% CP-VI-0054
% target_chapter: VI/03 — Coordinate Rings
% solution_provenance: FRESH_CANONICAL_AUTHORED_SOLUTION

\subsection*{CP-VI-0054: Coordinate ring of a graph}

Let $k$ be an infinite field and let $X=V(y-x^3)\subseteq\mathbb A_k^2$.
With $I(X)$ the ideal of polynomials vanishing at every $k$-point of $X$,
compute $k[X]=k[x,y]/I(X)$ and exhibit an explicit $k$-algebra
isomorphism $k[X]\cong k[t]$.

\medskip
\noindent\textbf{Solution.}

Every point of $X$ has the form $(a,a^3)$ with $a\in k$.
Let $F\in I(X)$. Division in $k[x][y]$ by the monic polynomial
$y-x^3$ gives
\[
F(x,y)=Q(x,y)(y-x^3)+R(x).
\]
For every $a\in k$, evaluation at $(a,a^3)$ yields $R(a)=0$.
Since $k$ is infinite, this implies $R=0$. Conversely $y-x^3$
vanishes on $X$, so $I(X)=(y-x^3)$.
Now define $\phi:k[x,y]\to k[t]$ by $x\mapsto t$ and $y\mapsto t^3$.
The map is surjective, and division by $y-x^3$ shows
$\ker\phi=(y-x^3)$. Hence
\[
k[X]=k[x,y]/I(X)\cong k[x,y]/(y-x^3)\cong k[t].
\]

\noindent\emph{Field-hypothesis check.} If instead $k=\mathbb F_q$ and
$X$ denotes its $k$-rational points, then $x^q-x$ also vanishes on
$X$ and is not in $(y-x^3)$. Thus the infinite-field hypothesis
cannot be dropped for the stated definition of $k[X]$.
